\subsection{Reconstruction in modular dimension at least three}

The following argument uses all finite generic covers of the small
etale topos.  We give the passage from the smaller Galois quotient
in the cited reconstruction theorem to the full Galois group, and
then recover the structure sheaf and the entire site.  Throughout
this subsection $k=\overline{\Q}$; field isomorphisms may act on $k$.

\begin{result}{E.T8}{Known, with full-group reduction proved}{The precise birational reconstruction used here}
Let $K$ and $L$ be function fields of transcendence degree at least
three over $k$.  Write $G_K=\operatorname{Gal}(K^{\mathrm{sep}}/K)$
and similarly for $L$.  The canonical map
\[
 \operatorname{Isom}(L,K)\longrightarrow
 \operatorname{Isom}_{\mathrm{cont}}(G_K,G_L)/\operatorname{Inn}(G_L)
                                                               \tag{E.T33}
\]
is a bijection.  For a field isomorphism $\varphi:L\to K$, its image
is $g\mapsto\widetilde\varphi^{-1}g\widetilde\varphi$, where
$\widetilde\varphi:L^{\mathrm{sep}}\to K^{\mathrm{sep}}$ is any
extension.  The ambiguity in this extension is exactly the indicated
inner ambiguity.  The construction is compatible with restriction
to every open subgroup and with its associated finite field extension.
\end{result}
\begin{proof}[Derivation, including the full-group comparison]
Fix a prime $\ell$.  Let $\Pi_K^c\to\Pi_K$ be the maximal pro-$\ell$
abelian-by-central quotient and its abelianization.  We use
\cite[Theorem 1.2(4), with Theorem 1.1]{Pop2018Divisorial} in its
actual form: the isomorphisms $\Pi_K\to\Pi_L$ which lift to
$\Pi_K^c\to\Pi_L^c$, modulo multiplication by a common
$\ell$-adic unit, correspond bijectively to field isomorphisms
$L\to K$.  The Kronecker dimension of $\overline\Q$ is one,
so its inequality is exactly $\operatorname{trdeg}_kK>2$.
The algebraically closed constant field is the unique maximal
algebraically closed subfield of the function field, and is
therefore preserved by the recovered isomorphism.  No specified
arithmetic action is added to the input.

Let $\alpha:G_K\to G_L$ be a continuous isomorphism.  The quoted
theorem first gives a unique $\varphi:L\to K$ from its induced
abelianized quotient.  For each open normal subgroup $H\triangleleft
G_K$, set
\[
 K_H=(K^{\mathrm{sep}})^H,\qquad
 L_{\alpha H}=(L^{\mathrm{sep}})^{\alpha H}.
\]
Both still have the required transcendence degree.  Applying the
same reconstruction to $\alpha|_H$ gives a unique
$\varphi_H:L_{\alpha H}\to K_H$.

We verify the compatibility which is needed to assemble these
maps.  For $g\in G_K$, conjugation by $g$ on $H$ and conjugation
by $\alpha(g)$ on $\alpha H$ commute with $\alpha|_H$.
Scalar multiplication on the abelianized quotients commutes with
these conjugations.  The uniqueness part of the field
reconstruction therefore gives
\[
 \varphi_H\circ\alpha(g)|_{L_{\alpha H}}
       =g|_{K_H}\circ\varphi_H.                            \tag{E.T34}
\]
Taking fixed fields shows that $\varphi_H$ carries $L$ onto $K$.
Its restriction to $L$ equals $\varphi$, for the following exact
reason.  The map $\Pi_{K_H}\to\Pi_K$ induced by $H\hookrightarrow
G_K$ has image of finite index.  On this image, the abelianized
map induced by $\varphi_H|_L$ agrees with that of $\alpha$, up to
the single $\ell$-adic scalar supplied for $\varphi_H$.
The group $\Pi_L$ is torsion-free: Kummer theory over the field
containing all roots of unity identifies it with a product of
copies of $\Z_\ell$, dual to the free abelian group $L^\times/k^\times$.
Two continuous $\Z_\ell$-linear maps to it agreeing on a subgroup
of finite index agree everywhere, by multiplying by that index.
The uniqueness assertion in the quoted theorem now proves
$\varphi_H|_L=\varphi$.

The identical argument for two nested open normal subgroups
$H_1\subset H_2$ proves
$\varphi_{H_1}|_{L_{\alpha H_2}}=\varphi_{H_2}$.
Their union is consequently a field isomorphism
\[
 \widetilde\varphi:L^{\mathrm{sep}}\xrightarrow{\sim}K^{\mathrm{sep}}.
\]
Equation (E.T34) on every finite Galois subextension says exactly
$\widetilde\varphi\alpha(g)=g\widetilde\varphi$ on the whole
separable closure.  This proves surjectivity in (E.T33).
Injectivity follows already from the smaller quotient theorem.
Applying the construction to an arbitrary open subgroup, and then
passing to its normal core, proves the final compatibility claim.

Here are the checked proof locations and the specialization scope
of the input.  Its reconstruction proceeds through the total
decomposition graph (Proposition 2.4), the divisorial lattice
(2.8), the quasi-arithmetic lattice and its functoriality (2.24),
the rational projections (2.27), and the recovery of divisorial
inertia from the group (3.3).  We read these proofs, the
specialization steps 2.17--2.21, and Jossen's complete appendix.
For the characteristic-zero fields used here, resolution permits
smooth projective models and normal-crossings complements; all
specializations in that application are taken on their common
good-reduction open.  The prime-to-$\ell$ comparison there is the
comparison used in 2.15 and 2.19.  The specialization argument
in 2.24 detects a non-torsion divisor class by reductions whose
order is divisible by $\ell$, and then restricts to a general
hyperplane; its residual transcendence degree is at least two.
This checks the exact dimension bound used in this proof.

The unrestricted reduced-special-fibre assertion in that paper's
Fact 2.11(1) must be restricted to the good-reduction locus of
Fact 2.13.  For example, at a valuation of $\overline\Q$ above
the prime $p$, the primitive projective conic
$X^2-pYZ=0$ has closure with special fibre $X^2=0$.
Its generic fibre is smooth, while that special fibre is a
doubled line.  The primitive-polynomial Gauss valuation shows
that its closure ideal is exactly $(X^2-pYZ)$, so no implicit
reduction changes this calculation.  Inverting $p$ removes
these specializations.  The cited reconstruction proof after
Fact 2.13 uses precisely the good-reduction open, and the
characteristic-zero application above uses only that form.
\end{proof}

% The two-dimensional extension is placed after E.T10 below.

\begin{result}{E.T9}{Proved}{Reconstruction of the full unringed topos in dimension at least three}
For normal separated varieties $X,Y$ over $k$ of dimension at least
three, the map
\[
 \operatorname{Isom}_{\mathrm{schemes}}(X,Y)
   \longrightarrow
 \operatorname{Isom}(X_{\mathrm{et}},Y_{\mathrm{et}})
                                                               \tag{E.T35}
\]
is a bijection, where the right side is taken modulo natural
isomorphism and the left side allows the induced automorphism
of $k$.  In particular, every unringed topos equivalence recovers
the structure sheaf and its action on every etale object.
\end{result}
\begin{proof}
We first reconstruct the map of locally ringed spaces.  Let
$F:X_{\mathrm{et}}\to Y_{\mathrm{et}}$ be an equivalence.
Subterminal objects are exactly the Zariski opens: a subsheaf
of the terminal object is represented by the union of the
images of its etale sections.  Consequently $F$ gives an
isomorphism of open-set lattices and a homeomorphism
$h:|X|\to|Y|$.  These spaces are sober.  Each irreducible
component, its generic point, and its codimension-one points
are recovered by this homeomorphism.  We work on one connected
component; a normal connected variety is integral.

Localizing along all nonempty Zariski opens recovers the generic
topos, namely the topos of continuous sets for $G_{k(X)}$.
The induced equivalence gives its continuous outer isomorphism
to $G_{k(Y)}$.  E.T8 gives a unique field isomorphism
\[
                \varphi:k(Y)\xrightarrow{\sim}k(X).          \tag{E.T22}
\]
We check that it has the valuations dictated by $h$.

Let $x$ be a codimension-one point of $X$ and $y=h(x)$.
Choose a geometric point over $x$ and the corresponding point
over $y$.  Every point of the small etale topos is induced by
a geometric point, by \cite[Lemma 59.29.12]{StacksEtalePoints};
its restriction to subterminals identifies its underlying
Zariski point.  For a finite generic Galois set $E$, the stalk at
$x$ of its generic direct image is
\[
                         E^{I_x}.                          \tag{E.T23}
\]
Indeed $\mathcal O_{X,x}$ is a discrete valuation ring.
The fraction field of its strict henselization is the fixed
field of the inertia subgroup $I_x\subset G_{k(X)}$ for the
chosen prolongation.  Sections of the finite generic cover
over that strict localization are exactly its invariant
elements.  Both the generic localization and this stalk
construction are intrinsic to the topos, so $F$ identifies
(E.T23) with its counterpart at $y$.
The fixed-point functor is represented, on finite continuous
sets, by the profinite homogeneous space $G_{k(X)}/I_x$.
Pro-Yoneda therefore identifies the transported inertia
subgroup with $I_y$, up to conjugacy.

Distinct divisorial valuations have distinct inertia classes
in the full generic group.  To see the needed assertion
directly, pass to the maximal pro-$\ell$ abelian quotient.
Kummer characters pair its inertia line with a rational
function $a$ by $v(a)\in\Z_\ell$.  Two distinct discrete
divisorial valuations admit, by weak approximation, an
element with valuation zero at one and nonzero at the other.
Their valuation kernels and hence their inertia lines differ.
Since the Galois action of (E.T22) is the one supplied by $F$,
we obtain equality of valuation rings
\[
                 \varphi(\mathcal O_{Y,y})=\mathcal O_{X,x}.
                                                               \tag{E.T24}
\]
This argument also shows that the scale of a chosen inertia
generator is immaterial: its full subgroup recovers the
valuation ring used in (E.T24).

For every open subset $V\subset Y$, normality gives
\[
 \mathcal O_Y(V)=\bigcap_{y\in V^{(1)}}\mathcal O_{Y,y}
                    \quad\hbox{inside }k(Y).                \tag{E.T25}
\]
The corresponding formula holds on $h^{-1}V$.  On an affine
open it is the normal-domain intersection theorem
\cite[Lemma 10.157.6]{StacksNormalIntersection}; on any open
it follows by covering with affine opens and comparing
inside the same fraction field.  Equations (E.T24)--(E.T25)
therefore give a sheaf isomorphism
$\mathcal O_Y\to h_*\mathcal O_X$ whose map at the generic
point is $\varphi$.  On stalks these are ring isomorphisms
between local rings and therefore preserve their unique
maximal ideals.  Thus it defines a scheme isomorphism $f:X\to Y$.
Its action on $k$ is the restriction of $\varphi$.

We finally prove that this recovers $F$ itself, including its
possible kernel on generic restriction.  The category of
quasicompact separated etale $X$-schemes is an intrinsic
generating subcategory of $X_{\mathrm{et}}$.  Here is the
precise recognition.  Every sheaf is represented by its
etale algebraic space: take the disjoint union of all its
etale sections and quotient by the open equivalence relation
where the sections agree.  Coherence makes this algebraic
space quasicompact and quasiseparated.  Its diagonal is
complemented as a subobject exactly when it is separated.
A separated quasicompact etale algebraic space over a scheme
is a scheme, by the quasi-finite separated form of Zariski's
Main Theorem.  Thus the coherent objects with complemented
diagonal are precisely this generating subcategory.

For such an object $U\to Y$, the slice equivalence identifies
$U_{\mathrm{et}}$ with the small etale topos of the
corresponding object $V\to X$.  Their generic fields are
finite extensions of $k(Y)$ and $k(X)$, respectively.
E.T8's open-subgroup compatibility gives the extension of
$\varphi$ between these fields.  Applying (E.T23)--(E.T25)
to the slice equivalence supplies a unique scheme
isomorphism $V\to U$ with that generic map, over $f$.
It is equivalently an isomorphism $V\simeq U\times_YX$.
These isomorphisms commute with every morphism in the
generating site: they agree on the generic fields, and
maps between separated integral schemes agreeing on a
dense open agree everywhere.  They also preserve all
covering families.  Hence they extend to a natural
isomorphism $F\simeq f_{\mathrm{et}}$ on the whole topos.
Uniqueness of $\varphi$ and (E.T25) proves injectivity in
(E.T35), completing the proof.
\end{proof}

\begin{result}{E.T10}{Proved}{The full unringed modular groups in dimension at least three}
For every stable type with $d=3g-3+n\geq3$, the canonical
map from semilinear algebraic stack autoequivalences of
$\mathcal M_{g,n,\overline\Q}$, modulo $2$-isomorphism, to
autoequivalences of its full unringed small etale topos,
modulo natural isomorphism, is an isomorphism.  Consequently
\[
 \Aut\bigl((\mathcal M_{g,n,\overline\Q})_{\mathrm{et}}\bigr)
 \cong
 \begin{cases}
 C_2\times G_{\Q},&(g,n)=(2,0),\\
 S_n\times G_{\Q},&\text{all other stable types with }d\geq3.
 \end{cases}                                                   \tag{E.T26}
\]
The maps from the factors are respectively the quadratic
twist of E.T4 or relabelling, and base change by an
automorphism of $\overline\Q$.  Thus this formula specifies
the comparison maps, as well as the abstract groups.
\end{result}
\begin{proof}
Put $\mathcal X=\mathcal M_{g,n,k}$.  We give the atlas
recognition and descent argument rather than adding a
structure sheaf to the topos as extra data.

We first record how stabilizers are seen by the topos.
For a closed geometric point $x$ of a separated
Deligne--Mumford stack of finite type over $k$, let $p_x^*$
be its stalk functor.  There is a canonical isomorphism
\[
                   I_x\xrightarrow{\sim}\Aut(p_x^*),         \tag{E.T27}
\]
where $I_x$ is the finite automorphism group of the object
$x$.  Choose a representable etale scheme atlas $W$ and a
lift $w$ of $x$.  An automorphism of $p_x^*$ acts on its
set of lifts to $W$.  The open subschemes of $W$ force the
underlying closed point of $w$ to be fixed: for any other
closed point there is an open containing $w$ and excluding
that point.  The lifts with this same underlying point form
an $I_x$-torsor, because the atlas has trivial stabilizer.
After composing with the unique element of $I_x$ sending
the chosen lift to its image, the stalk automorphism fixes
that lift.  It is then a stalk automorphism of $W_{\rm et}$
at $w$.  Such an automorphism is the identity: on every
quasicompact etale $W$-scheme its finite set of lifts over the closed
$k$-point is separated by open neighbourhoods.  This proves
surjectivity in (E.T27).  The torsor action proves injectivity.
The same proof after an etale chart shows that all points of
the stack topos arise from geometric points, and that the
closed ones are recognized by their restriction to the
subterminal open-set lattice.

Choose a finite etale surjective level atlas
$U\to\mathcal X$ with $U$ a smooth quasiprojective scheme.
For genus zero we take $U=\mathcal X$; in positive genus
one may take full symplectic level $3$, with the labels
retained.  The level structure is a torsor under a finite
symplectic group.  It kills curve automorphisms: the
principal level-$3$ subgroup has no nontrivial finite-order
integral matrix, and a finite-order curve automorphism
acting trivially on integral first homology is the identity.
For genus at least two the latter follows from the
quotient curve and Riemann--Hurwitz, since invariant
holomorphic differentials would give the quotient the
same genus; in genus one the marked origin removes
translations.  The resulting moduli space is the usual
fine level moduli scheme.  Its finite etale atlas and the
quasiprojective coarse moduli scheme are the same finite
level models used in the completion construction of E.2.

Let $F$ be an equivalence of $\mathcal X_{\rm et}$ and
write $V=F^*U$.  A finite locally constant object on a
Deligne--Mumford etale site is a representable finite etale
cover: on an etale scheme atlas it is a finite disjoint
union of finite etale schemes, and the sheaf descent
isomorphisms give effective finite etale descent.  Hence
$V$ is a finite etale Deligne--Mumford stack over
$\mathcal X$, and the slice topos $V_{\rm et}$ is
equivalent to $U_{\rm et}$.  Every closed point of the
latter has trivial stalk automorphism group.  Formula
(E.T27) gives trivial stabilizers at every closed point
of $V$.  Its inertia is finite and unramified; a nonempty
complement of its identity section would have a closed
point, since $V$ is of finite type over $k$.  Its inertia
is therefore the identity section, and $V$ is an
algebraic space.

We verify that this space is a scheme.  Let $M$ be the
coarse moduli scheme of $\mathcal X$.  The map
$\mathcal X\to M$ is proper and quasi-finite; its
geometric fibres are the classifying stacks of the
finite curve automorphism groups.  It follows that
$V\to M$ is proper and quasi-finite.  A proper
quasi-finite algebraic space over a scheme is finite
and hence a scheme; alternatively the separated
quasi-finite assertion of
\cite[Proposition 67.50.2]{StacksSeparatedQuasiFinite}
first makes it a scheme, followed by the scheme form
of the same finiteness statement.  Thus $V$ is a smooth
scheme of pure dimension $d$.

Set $R=U\times_{\mathcal X}U$.  It is a scheme finite
etale over $U$, because $U\to\mathcal X$ is so.  Its
image under $F^*$ is
$R'=V\times_{\mathcal X}V$, again a scheme finite
etale over $V$.  The slice equivalences and E.T9 give
unique scheme isomorphisms
\[
                    V\longrightarrow U,
          \qquad R'\longrightarrow R.                       \tag{E.T28}
\]
E.T9 applies to each connected component, whose dimension
is $d\geq3$.  Its construction is natural for etale maps:
on the generic fields this is E.T8's compatibility with
finite extensions, and equality on the dense generic
locus extends over a normal separated scheme.  Thus
(E.T28) commutes with source, target, identity, inverse,
and composition.  For composition apply the same
argument to $R\times_U R$ and its image; its maps are
etale as well.  All field isomorphisms in the same
connected component of this groupoid have the same
restriction to $k$.  The stack $\mathcal X$ is connected,
so the relations connect the components of the atlas,
and there is one common automorphism of $k$.

The isomorphism of etale groupoids (E.T28) descends to
a semilinear stack autoequivalence $f$ of
$[U/R]=\mathcal X$.  On each atlas slice it induces
the original topos equivalence by E.T9, and the
agreement on the overlap slices gives
$f_{\rm et}\simeq F$ on the full topos by descent.
Conversely, a natural isomorphism of such topos
equivalences yields the unique isomorphisms on their
scheme-atlas groupoids from E.T9, and hence a
$2$-isomorphism of the stack maps.  This proves the
claimed bijective comparison, including its kernel.

Restriction to constants gives an exact sequence with
kernel the $k$-linear algebraic group $A_{g,n}$ of E.T4
and quotient $G_{\Q}$.  The canonical model of the
moduli stack over $\Q$ splits it.  Relabellings are
defined over $\Q$, so this splitting centralizes $S_n$.
In type $(2,0)$ the kernel is $C_2$, which has trivial
automorphism group, so the split extension is again
the direct product.  Substituting E.T4 gives (E.T26).
\end{proof}
